\documentclass[10pt,a4paper]{amsart}
\usepackage[margin=19mm]{geometry}
\usepackage{amsmath,amssymb,mathtools}
\usepackage[scr=rsfs,cal=euler]{mathalfa}
\usepackage{xfrac}
\usepackage[unicode,hidelinks,bookmarks=false]{hyperref}
\newcommand{\ie}{i.e.\ }
\newcommand{\cf}{cf.\ }
\newcommand{\resp}{respectively\ }
\newcommand{\Subsection}{\subsection}
\providecommand{\nobreakdash}{\nobreak\hbox{-}}
\setlength{\emergencystretch}{2em}
\multlinegap0pt
\usepackage{bm}
\usepackage{bbm}
\usepackage{dsfont}
\usepackage{xspace}
\usepackage{cancel}
\usepackage{mathtools}
\usepackage{amsthm}
\makeatletter
\@ifundefined{theorem}{\newtheorem{theorem}{Theorem}}{}
\@ifundefined{definition}{\newtheorem{definition}[theorem]{Definition}}{}
\@ifundefined{lemma}{\newtheorem{lemma}[theorem]{Lemma}}{}
\makeatother
\makeatletter
\newcommand{\N}{\mathbb{N}}
\newcommand{\R}{\mathbb{R}}
\expandafter\ifx\csname abstract\endcsname\relax
\renewenvironment{abstract}{
\begin{center}
{\bfseries \large\abstractname\vspace{\z@}}
\end{center}
\quotation
}
\fi
\renewcommand{\epsilon}{\varepsilon}
\renewcommand{\geq}{\geqslant}
\renewcommand{\leq}{\leqslant}
\renewcommand{\mod}{\operatorname{mod}}
\newcommand{\oo}{\infty}
\renewcommand{\subset}{\subseteq}
\makeatother
\title[Weighted averages and applications to sets of multiple recur]{Weighted averages and applications to sets of multiple recurrence}
\author{Vitaly Bergelson, Michael Reilly}
\date{Source: arXiv:2609.16554v1 (CC BY 4.0)}
\begin{document}
\maketitle

\noindent\emph{Source and licence.} Weighted averages and applications to sets of multiple recurrence. Version arXiv:2609.16554v1 (CC BY 4.0), distributed under the Creative Commons Attribution 4.0 International licence, as recorded in the arXiv licence metadata for this exact version and shown on the abstract page https://arxiv.org/abs/2609.16554v1. Full rights notice: licenses/arxiv-2609.16554-CC-BY-4.0.txt.\par\smallskip

\noindent\textit{Editorial scope.} Counted unit: the Lemma whose source label is \texttt{lem:derivatives\_mod\_1} in the article (source0.tex lines 1074--1077, source label \texttt{lem:derivatives\_mod\_1}), delivered together with the article's own complete original proof of it (source0.tex lines 1078--1133, one \texttt{proof} environment of 56 lines). No mathematics is written, repaired or re-derived: statement and proof are the article's own text line for line. Required context retained uncounted: def:W\_thick (lines 570--575); lem:technical\_1 (lines 1041--1043); lem:shift\_inequality (lines 1052--1057). Every definition, display and label of those lines is retained unchanged. Omitted from this package and not redistributed: 1--569, 576--1040, 1044--1051, 1058--1073, 1134--2501. Nothing inside a retained excerpt is omitted or altered: every retained body is delivered exactly as the article's lines are written, so no omission receipt of any kind accompanies this package. Citations: the retained excerpts carry no citation command at all, so no bibliography entry of the article is transcribed and no reference is redistributed. Reference adaptation: none was needed and none was made; every reference inside the delivered unit resolves inside it, so no unresolved reference or citation is produced. Printed slips kept literal: no printed word, symbol, hypothesis or number of the counted statement or of its proof was altered. Layout and class: the article's own class is \texttt{article} with packages including babel, inputenc, geometry, tocloft, amsfonts, mathrsfs, bbm, latexsym, amssymb, mathtools, relsize, lipsum, todonotes, enumitem; the wrapper is the lane's pinned \texttt{amsart} editorial wrapper at 10pt on A4, which restates the theorem-like environments the retained text needs and adds the article's own macro definitions that the retained text needs (\texttt{\textbackslash N}, \texttt{\textbackslash R}, \texttt{\textbackslash epsilon}, \texttt{\textbackslash geq}, \texttt{\textbackslash leq}, \texttt{\textbackslash mod}, \texttt{\textbackslash oo}, \texttt{\textbackslash subset}), with extra wrapper packages (bm, bbm, dsfont, xspace, cancel, mathtools). The delivered numbering is the unit's own. Assets: no retained span contains a figure environment, a graphics-inclusion command, a PDF-page inclusion, a table or a TikZ picture; no image file is copied into this package, the package declares no asset dependency and no image file is redistributed with it. Licence: version arXiv:2609.16554v1 of this article is distributed under the Creative Commons Attribution 4.0 International licence, as recorded in the arXiv licence metadata for this exact version and shown on the abstract page https://arxiv.org/abs/2609.16554v1. A full rights notice is delivered at licenses/arxiv-2609.16554-CC-BY-4.0.txt. Characters: every retained excerpt is pure ASCII, so the delivered engine, XeLaTeX with its default Latin Modern text font, reports no missing character.
\begin{definition}\label{def:W_thick}
    Let $W$ be a function with $1\prec W(x)\preceq x$. We say that $S$ is $W$\emph{-thick} if 
    \begin{equation}
        \limsup_{W(N)-W(M)\to\oo}\frac{1}{W(N)-W(M)}\sum_{n=M}^N\Delta W(n)1_{S}(n)=1.
    \end{equation}
\end{definition}

\begin{lemma}\label{lem:technical_1}
    Let $f: \N\rightarrow \R$ be an increasing function, let $A<B<C$ be elements of $(0,1)$ with $B<C-A$, and let $X<Y$ be natural numbers with $Y-X> \frac{2}{A}$. Suppose that $(\Delta f(n)\text{ mod } 1)\in (A,B)$ for all $n\in \{X,\dots, Y\}$. Then there exist natural numbers $Z,W$  such that $ (f(n)\text{ mod } 1)\in (A,C)$ for all $n\in \{Z,\dots, W\}$ with $X\leq Z<W$, $Z< X+\frac{1}{A}$, $W-Z\geq \frac{C-A}{B}-2$.
\end{lemma}

\begin{lemma}\label{lem:shift_inequality}
    Let $g$ be a Hardy function with $\lim_{x\to\oo}g(x)=\oo$ and $\lim_{x\to\oo}g'(x)=0$. Fix any $\epsilon\in (0,1)$. Then
\begin{equation}\label{eq:shift_inequality}
        g'\left(x+\frac{1}{g'(x)^{\epsilon}}\right)=g'(x) \cdot (1+o_{x\to\oo}(1)). 
    \end{equation}
\end{lemma}

\begin{lemma}\label{lem:derivatives_mod_1}
 Let $f$ be a Hardy function which increases to $\oo$. Let $d\in \N$ with $d>1$ and suppose that $x^{d-1}\prec f(x) \prec x^{d}$. 
 Then for each $\epsilon>0$, there are arbitrarily large values of $N\in \N$ with $\Delta^{i}f(N)\text{ mod } 1\in (0,\epsilon\cdot (\Delta^{d}f(N))^{i/d})$ for all $i\in\{0,1,\dots, d\}$. Additionally, $f(n)\text{ mod } 1\in (0,\epsilon)$ for all $n\in [N,N+K(N)]$, where $K$ is a function that satisfies $\lim_{N\to\oo}\frac{K(N)}{(\Delta^{d} f(N))^{-1/d}}=\oo$, and so it follows that the set $\{n\in \N: f(n)\mod 1\in (0,\epsilon)\}$ is $W$-thick (Definition \ref{def:W_thick}) for $W(N) = (\Delta^{d} f(N))^{-1/d}$.
\end{lemma}

\begin{proof}

We begin by noting that $\Delta^df$ decreases to $0$. Let $\epsilon\in (0,1)$ and let $s:\N\rightarrow \N$ be a function which increases to $\oo$ such that
\begin{equation}\label{eq:s_growth}
\lim_{N\to\oo}s(N)\cdot \Delta^df(N)=0\quad \text{ and }\quad \lim_{N\to\oo}s(N)\cdot (\Delta^df(N))^{\frac{d-1}{2d-1}}=\oo
\end{equation}
(for example, one could take $s(n) = (\Delta^d f(n))^{-1+1/(2d-1)}$).
For $n\in \N$, put $A(n) =s(n)\cdot \Delta^df(n)$ and for each $i\in \{1,\dots, d\}$ put 
\begin{equation}
   B_{i}(n) = \epsilon\cdot \left({\Delta^{d}f(n)}\cdot s(n)\right)^{1-i/d}\quad \text{ and }\quad I_i(n) = \left[A(n),B_i(n)\right].
\end{equation}
Observe that 
\begin{align*}
    \frac{A(n)}{B(n)} = (s(n)\cdot \Delta^df(n))^{1/d}=o_{n\to\oo}(1).
\end{align*}
It follows that the length of $I_i(n)$ is $B_i(n)\cdot (1+o_{n\to\oo}(1)) $. Next, pick an arbitrarily large value of $N_1\in \N$ such that
\[
(\Delta^{d-1}f(N_1-1) \text{ mod } 1) \not\in I_{1}(N_1)\text{ and }(\Delta^{d-1}f(N_1) \text{ mod } 1) \in I_{1}(N_1).
\]
Such a value of $N_1$ exists because $\Delta^{d-1}f$ increases to infinity, $\Delta^{d}f$ decreases to $0$, and the interval $I_{1}(N_1)$ has length much larger than $\Delta^{d}f(N_1)$ when $N_1$ is large enough. From now on, we put $A = A(N_1)$, $B_i = B_i(N_1)$, and $I_i = I_i(N_1)$ for $i\in \{1,\dots d\}$. Note that for each $i\in \{1,\dots, d\}$ we have that $B_i-A>B_{i-1}$ so long as $N_1$ is large enough.  


For $n\geq N_1$, the sequence $(\Delta^{d-1}f(n) \text{ mod } 1)_{n\in \N}$ takes steps of size $\Delta^{d}f(n)\leq \Delta^{d}f(N_1)$ and so it follows that $(\Delta^{d-1}f(n) \text{ mod } 1) \in I_{1}$ for $n\in \{N_1,\dots, N_1+K_1\}$, where $K_1$ satisfies
\begin{equation}\label{eq:K_1_big}
K_1\geq \frac{ |I_{1}|}{\Delta^{d}f(N_1)}-2 = \epsilon\cdot {\Delta^{d}f(N_1)^{-1/d}}\cdot s(N_1)^{1-1/d} \geq 2/ A
\end{equation}
when $N_1$ is large enough, since 
\[
\frac{B_1 A}{\Delta^{d}f(N_1)}= (\Delta^{d}f(N_1))^{1-1/d}\cdot s(N_1)^{2-1/d} = ((\Delta^{d}f(N_1))^{(d-1)/(2d-1)}\cdot s(N_1))^{2-1/d}\to \oo
\]
as $N_1\to\oo$.


\textbf{Claim:} There exist natural numbers, $N_1\leq \dots \leq N_{d}$ and $K_1,\dots, K_{d}$ with $N_{i}\leq N_1+\frac{i-1}{A}$ and $K_1\geq 2/A$, such that $(\Delta^{d-i}f(n)\text{ mod } 1)\in  I_{i}$ for all $n\in \{N_i,\dots, N_i+K_i\}$ and all $i\in\{1,\dots d\}$.


We prove this claim by induction on $i$. We have already shown the base case $i=1$, and now we show the induction step.

Suppose that $i\in\{1,\dots, d-1\}$ and that $N_i$ and $K_i$ are integers for which we have $(\Delta^{d-i}f(n) \text{ mod } 1) \in I_{i}$ for $n\in \{N_i,\dots, N_i+K_i\}$, where $N_i\leq N_{1}+\frac{i-1}{A}$ and $K_i\geq 2/A$. By Lemma \ref{lem:technical_1}, there exist $N_{i+1},K_{i+1}\in \N$ such that $(\Delta^{d-(i+1)}f(n) \text{ mod } 1) \in I_{i+1}$ for all $n\in\{N_{i+1},\dots, N_{i+1}+K_{i+1}\}$, where $N_i\leq N_{i+1}\leq N_{i}+\frac{1}{A}< N_{1}+\frac{i}{A}$ and 
\begin{equation}\label{eq:length_of_interval}
K_{i+1}\geq \frac{ B_{i+1}-A}{B_i}-2 =  \left(\Delta^{d}f(N_1)\cdot s(N_1)\right)^{-1/d} \cdot (1+o_{N_1}(1))\geq 2/A
\end{equation}
when $N_1$ is large enough. This completes the induction step and the proof of the claim. 


Next, since $N_{d}\leq N_1+\frac{d-1}{A}$ we have $N_{d}\leq N_1+K_1$ by (\ref{eq:K_1_big}). Thus, taking $N=N_{d}$, we have that $N$ lies in each of the intervals $\{N_i,\dots, N_i+K_i\}$ for $i\in \{1,\dots, d\}$. So,
\[
(\Delta^{d-i}f(N)\text{ mod } 1)\in I_i\subset(0,\epsilon\cdot (\Delta^{d}f(N_1))^{\frac{d-i}{d}}) \text{ for all }i\in\{1,\dots d\}.
\]

Recall that $\Delta^{d}f(N_1)\sim  \Delta^{d}f\left( N_1+\frac{d-2}{A}\right)=\Delta^{d}f\left( N_1+\frac{d-2}{(\Delta^{d}f(N_1))^{\eta}}\right)$ from Lemma \ref{lem:shift_inequality}. Replacing $\epsilon$ with $\epsilon/2$ if necessary, we have
\[
(\Delta^{i}f(N)\text{ mod } 1)\in (0,\epsilon\cdot (\Delta^{d}f(N))^{i/d})
\]
for all $i\in\{1,\dots d\}$. This shows the first claim in the statement of the Theorem. The second claim is immediate by taking $K = K_{d}$ and applying  Lemma \ref{lem:shift_inequality} as above.
\end{proof}

\end{document}
